\subsubsection*{Solution}
The requested quantity does not presently have a rigorously established exact value. In fact, the commonly presumed result $f(\gamma)=1-\gamma$ is false; the exact relative-entropy contraction coefficient of the amplitude-damping channel remains an optimization problem. The relevant literature gives rigorous bounds and a numerically supported conjecture, but not a proof of that conjecture. See \href{\detokenize{https://arxiv.org/abs/1508.03551}}{Hiai and Ruskai} and Section 6 of \href{\detokenize{https://arxiv.org/abs/2411.19890}}{Belzig-Gao-Smith-Wu}.

\paragraph*{Step-by-Step Derivation}

Set

\[
q=1-\gamma.
\]

\subparagraph*{1. Bloch-sphere action}

Write a qubit state as

\[
\rho=\frac12\left(I+x\sigma_x+y\sigma_y+z\sigma_z\right).
\]

The amplitude-damping channel acts affinely as

\[
(x,y,z)\longmapsto \left(\sqrt q\,x,\sqrt q\,y,\gamma+qz\right).
\]

Therefore, for the difference of two states,

\[
(\Delta x,\Delta y,\Delta z) \longmapsto \left(\sqrt q\,\Delta x,\sqrt q\,\Delta y,q\,\Delta z\right).
\]

For qubit states, the trace norm of their difference is the Euclidean norm of the Bloch-vector difference. Hence the trace-distance contraction coefficient is

\[
\eta_{\mathrm{tr}}(\mathcal A_\gamma) = \max_{\Delta\mathbf r\ne0} \frac{\sqrt{q(\Delta x^2+\Delta y^2)+q^2\Delta z^2}} {\sqrt{\Delta x^2+\Delta y^2+\Delta z^2}} =\sqrt q.
\]

The general comparison theorem for quantum channels gives

\[
\eta_{\mathrm{tr}}(\mathcal N)^2 \leq \eta_D(\mathcal N) \leq \eta_{\mathrm{tr}}(\mathcal N).
\]

Thus,

\[
1-\gamma\leq f(\gamma)\leq\sqrt{1-\gamma}.
\]

Consequently,

\[
\frac78+\frac34+\frac12 \leq f\left(\frac18\right)+f\left(\frac14\right)+f\left(\frac12\right) \leq \sqrt{\frac78}+\sqrt{\frac34}+\sqrt{\frac12},
\]

or

\[
{ \frac{17}{8} \leq f\left(\frac18\right)+f\left(\frac14\right)+f\left(\frac12\right) \leq \frac{\sqrt{14}+2\sqrt3+2\sqrt2}{4} }.
\]

Numerically,

\[
2.125 \leq f\left(\frac18\right)+f\left(\frac14\right)+f\left(\frac12\right) \leq 2.50854653166447.
\]

\subparagraph*{2. Why $f(\gamma)\neq1-\gamma$}

Consider

\[
\sigma_p=
\begin{pmatrix}
1-p&0\\
0&p
\end{pmatrix},
\qquad
\rho_{p,\varepsilon}
=
\sigma_p+\varepsilon\sigma_x,
\]

where $\varepsilon^2<p(1-p)$. Define

\[
h(p)= \frac{\log\!\left(\frac{1-p}{p}\right)}{1-2p}, \qquad h\left(\frac12\right)=2.
\]

The second-order expansion of relative entropy is

\[
D(\rho_{p,\varepsilon}\Vert\sigma_p) = \varepsilon^2h(p)+O(\varepsilon^3).
\]

The channel sends

\[
\sigma_p\longmapsto
\begin{pmatrix}
1-qp&0\\
0&qp
\end{pmatrix},
\qquad
\varepsilon\sigma_x\longmapsto
\sqrt q\,\varepsilon\sigma_x.
\]

Therefore,

\[
D\!\left( \mathcal A_\gamma(\rho_{p,\varepsilon}) \middle\| \mathcal A_\gamma(\sigma_p) \right) = q\varepsilon^2h(qp)+O(\varepsilon^3),
\]

so

\[
f(\gamma)\geq q\frac{h(qp)}{h(p)}.
\]

For example, choosing $p=\tfrac12$ and $\gamma=\tfrac12$ gives

\[
f\left(\frac12\right) \geq \frac12\,\frac{h(1/4)}{h(1/2)} = \frac{\log3}{2} = 0.549306144334055 > \frac12.
\]

Thus the tempting answer $17/8$, based on $f(\gamma)=1-\gamma$, cannot be correct.

\subparagraph*{3. Conjectured numerical result}

The numerically supported-but unproved-formula reported in the literature is

\[
f(\gamma)\stackrel{?}{=} \max_{0<p<1} (1-\gamma) \frac{h((1-\gamma)p)}{h(p)}.
\]

Numerical maximization gives

\[
\begin{aligned}
f\left(\frac18\right)&\stackrel{?}{=}0.905385896325533,\\
f\left(\frac14\right)&\stackrel{?}{=}0.807203927129647,\\
f\left(\frac12\right)&\stackrel{?}{=}0.596500342973557.
\end{aligned}
\]

Their sum is

\[
2.30909016642874.
\]
